\scanpage{53}
% manuscript line 1
hilbertienne déduite de celle de $\bar E\overset{(2)}{\otimes}F$ par transport
de structure, est aussi noté $L^{(2)}(E,F)$, \inserted{ou $L^2(E)$ \deleted{si} si $F=E$.} La norme de \deleted{l'op.}
\deleted{\unsure{la} \illegible{1}} l'élément $u\in L^2(E,F)$ est notée $\|u\|_2$.

% manuscript line 4
\[
\text{On a }\deleted{\text{\unsure{Rappelons} donc}}\quad\boxed{\|u\|_\infty\leq\|u\|_2} \tag*{(16)}
\]

% manuscript line 5
Les op. de Hilbert-Schmidt de $E$ dans $F$
sont des applications compactes (car limites dans $L(E,F)$ d'él. de rang fini. Ce sont
\deleted{\unsure{tout de suite}} \deleted{ceux dont} \inserted{aussi exactement les op. dont} les matrices $(u_{ij})$ par rapport à des bases
orthonormales $(e_i),(f_j)$, \deleted{sont} de carré sommable
(matrices dites aussi de H.-S.) \inserted{et on a \inserted{alors} $\|u\|_2=(\sum|u_{ij}|^2)^{1/2}$.} Par la suite, nous
identifions $\bar E\overset{(2)}{\otimes} F$ et $L^2(E,F)$.

% manuscript line 7a
% slanted
\margin{arranger \deleted{mettre le (Th. 2)} \uline{Une proposition}}

% manuscript line 11
\uline{\unsure{Théorème 2}} Soient $E$ et $F$ deux espaces de Hilbert,
soit $u$ une \deleted{op. de H.-S.} \inserted{application linéaire continue} de $E$ dans $F$
\deleted{et si $G$ est un \unsure{autre esp.} de Hilbert}

% manuscript line 14
1) Si $u$ est un op. de H.-S., il en est de $=$
de $Au$ (resp. $uA$) pour \deleted{\unsure{toute}} application linéaire
continue \deleted{de} \inserted{de $F$ dans} un espace de Hilbert $G$ \deleted{dans} (resp. de
$G$ dans $E$) \,; et on a

% manuscript line 18
\[
\deleted{(17)}\quad\boxed{\|Au\|_2\leq\|A\|\,\|u\|_2\quad
\text{(resp. }\|uA\|_2\leq\|A\|\,\|u\|_2\text{)}} \tag*{(17)}
\]

% manuscript line 19
2) Pour que $u$ soit un op. de Hilbert-Schmidt, il
faut et il suffit que $u^*$ soit un op. de
H.-S. (de $F$ dans $E$). \inserted{et alors $\boxed{\|u\|_2=\|u^*\|_2}$. Plus précisément,}
\deleted{On} a, pour $u,v\in L^{(2)}(E,F)$

% manuscript line 23
\[
\boxed{(u,v)=(v^*,u^*)} \tag*{(18)}
\]

% manuscript line 24
% slanted
\margin{\uline{Th. 3} (Caractérisation spectrale)}

% manuscript line 24
3) Soit $h=\sqrt{u^*u}\,\deleted{\unsure{\in} E}$ \inserted{(op. hermitien positif dans $E$)}. Pour que $u$ soit un
H.-S., il faut et suffit \deleted{\unsure{que} $h$ \unsure{soit} \illegible{3}} \deleted{que la suite \illegible{2} $\rho_i$,}
\deleted{\unsure{des valeurs propres} (\unsure{i|+}) de $h$ soit de carré sommable,}
\deleted{et alors} que l'on ait $\sum(\mathrm{v.s.})^2<+\infty$. Alors

% manuscript line 28
\[
\boxed{\|u\|_2=\|h\|_2=(\sum\rho_i^2)^{1/2}} \tag*{(19)}
\]

% manuscript line 29
\deleted{\unsure{parmi les valeurs propres de} $\sqrt{u^*u}$}

% manuscript line 29a
% slanted
\margin{\uline{Corollaire 1} $u$ H.-S. $\Longleftrightarrow |u|$ H.-S. alors $\|u\|_2=\||u|\|_2$}

% manuscript line 30
\uline{Corollaire 2} $u$ H.-S. $\Longleftrightarrow u^*u$ Fredholm, et on a

% manuscript line 31
\[
\|u\|_2=\sqrt{\operatorname{Tr}\deleted{\unsure{(u)}}u^*u}=\sqrt{\|u^*u\|_1}
\]

% manuscript line 32
% archive sheet number: 54
